\subsection{Triangle Inequality} 
If three points are \emph{not} \makeblank[1in], they form a triangle, which gives us the name of the following property of the Euclidean plane.
\begin{displaybox}[Triangle Inequality]
Given any three points $A$, $B$, and $C$:
\[
AC\leq AB+BC
\]
with equality if and only if the three points are colinear with $B$ on $\overline{AC}$.
\end{displaybox}
We sometimes phrase this as ``\makeblank[3.5in].''
\begin{Note}
This is primarily interesting when $\overline{AC}$ is the longest of the three segments.
\end{Note}
To practice our understanding of this inequality, we may ask whether a list of three numbers are possible \emph{sides of a triangle}, that is, possible lengths of three segments $AB$, $BC$, and $AC$.
\begin{Example}
Which of these are possible lists of lengths of the sides of a triangle?
\begin{enumerate}
\item $AB=5$, $BC=6$, $AC=10$: \makeblank, so \makeblanksmall
\item $AB=1$, $BC=6$, $AC=3$: \makeblank, so \makeblanksmall
\begin{Note}
This \makeblank[1.5in] for any three \makeblank[1in] points---this is \makeblank[1.5in].
\end{Note}
\item $AB=5$, $BC=7$, $AC=8$: \makeblank, so \makeblanksmall
\item $AB=2$, $BC=4$, $AC=6$: \makeblank, so \makeblanksmall
\begin{Note}
This is \makeblanksmall a triangle because the points are \makeblank[1in]. We might call this a ``degenerate triangle.''
\end{Note}
\end{enumerate}

\end{Example}